\section{Extending ClimateSOM: Detecting Distributional Shifts in Precipitation Co-variability}\label{sec:agu}

In this section, I continue with the \emph{ensemble variability} facet, sharpening it from visual exploration to a specific statistical question.
%
Climate-change impacts are typically assessed through mean-state changes at individual locations, yet the spatial co-variability \emph{between} regions may reorganize fundamentally under warming in ways that marginal trends alone cannot reveal.
%
Detecting such reorganization demands an explicit comparison of the distributions of spatial realizations across scenarios, something neither mean-state analysis nor EOF-based decompositions are designed to directly support.
%
To address this, I extend the ClimateSOM framework with statistical significance testing capabilities, to determine whether the \emph{distribution} of conditioned spatial precipitation patterns has shifted with statistical significance.
%
% Discovering such distributional shifts is difficult with standard analyses and here we demonstrate how ClimateSOM can be used to detect such changes under climate change.
%
I demonstrate this capability on the co-variability between California's Sacramento Bay-Delta Watershed and U.S. West Coast precipitation under LOCA2-downscaled CMIP6 projections, where I find a strengthening departure from historical relationships that I term ``decoherence,'' and the remainder of this section develops the methodology and findings in detail.
%
Although this extension remains a work in progress, I consider it a true litmus test of ClimateSOM's utility. 
%
Rather than a contrived demonstration, it applies the workflow to a climatologically substantive question, whether and how regional precipitation co-variability reorganizes under warming, and stakes its value on whether it can deliver insight that the standard analysis methods of the domain, such as mean-state trend analysis and EOF-based decomposition would struggle. 
%
It is in this confrontation with a real problem, and with the established tools climate scientists already trust, that the workflow's practical worth is genuinely put to the test.

%% Suggested section heads:
\subsection{Introduction}
\label{sec:introduction}
Climate change impacts are typically assessed through mean-state changes at individual locations or regions.
%
While such analyses provide important insights (e.g., projected wetting winters and drying springs in California~\cite{pierce2018climate,polade2017precipitation}), they overlook a critical dimension: spatial co-variability between regions may reorganize fundamentally under warming.
%
Even when marginal trends appear modest for one region, a shift in the co-variability between two regions can carry consequences that marginal trends alone do not reveal.
%
Detecting such reorganization requires an explicit comparison of distributions of spatial realizations across climate scenarios, as temporal or spatial means and correlations alone can mask compositional shifts in realized patterns.

Beyond mean-state analysis, Empirical Orthogonal Function (EOF) analysis is widely used to identify dominant modes of variability~\cite{hannachi2007empirical}, including covariability between climate modes~\cite{boschat2026covariability}. However, EOF-based approaches rely on linear combinations of orthogonal spatial patterns and may not capture nonlinear relationships or distributional shifts in realized patterns.
%
While extensions of EOFs, such as rotated EOFs~\cite{lian_evaluation_2012}, have been developed to address some of these limitations, they still primarily focus on identifying dominant modes of variability and do not directly support the distributional view that we pursue.

The value of viewing climate realizations from a distributional perspective has been shared and appreciated by past authors.
%
Pendergrass et al. examined the effects of a warmer climate considering the distribution of spatially-aggregated change in the interannual standard deviation of seasonal-mean precipitation, globally~\cite{pendergrass_precipitation_2017}.
%
Zhang et al. argues that warming produces a distributional shift in precipitation, with wider variability and more extremes~\cite{zhang_increasing_2021}.
%
Crucially, they highlight the unreliability of examining mean states alone; both because changes in mean states alone cannot inform changes in variability and because change in variability may, in fact, be more difficult to adapt to than changes in mean states.
%
A more recent work reinforced this idea by showing that day-to-day variability changes can outweigh mean-state shifts in driving precipitation extremes across many regions~\cite{nordling2025climate}.
%
However, these past works fall short of considering the distribution of the spatial realizations themselves, which we argue is critical to better understanding the impacts of climate change on regional precipitation patterns and their co-variability.

% However, we contend such analysis has been largely absent from the literature.
%
%
% framework enables rigorous hypothesis testing of whether the distribution of spatial precipitation patterns conditioned on regional climate states has shifted between historical and future periods.
To address this gap, we extend the ClimateSOM framework, building on prior applications of self-organizing maps to climate pattern analysis \cite{hewitson_self-organizing_2002, gibson_use_2017}. ClimateSOM discretizes the space of spatial climate patterns into a nonlinear, topology-preserving partition and thereby admits the distributional comparisons that EOF-based decompositions do not support. We extend this framework with a permutation-based hypothesis testing procedure utilizing Earth Mover's Distance to quantify distributional differences in spatial realizations.
% %
% Prior work has examined precipitation linkages and regional co-variability through correlation-based or mean-state analysis~\cite{}, but these approaches cannot formally test whether the composition of realized spatial patterns has shifted.
%
% Our framework enables such tests, 
% , and the stationarity assumptions underlying water resource planning.
We further develop a within-year coherence test and a clustering-based decomposition that attributes distributional shifts to compositional reorganization among identifiable pattern clusters.

We demonstrate this capability on a critical test case concerning associations between California's Sacramento Bay-Delta Watershed and U.S. West Coast precipitation.
%
The Bay-Delta is a critical hydrological hub in a persistently drought-prone region~\cite{cloern2011projected} and a focus of climate-change impact studies~\cite{rojas-aguirre_evaluating_2025}.
%
Precipitation on the U.S. West Coast has also been widely studied from its North-South gradient~\cite{dettinger1998north}, the role of atmospheric rivers in shaping seasonal totals and extremes \cite{gershunov_assessing_2017, gershunov_precipitation_2019}, and future precipitation projections in the Pacific Northwest~\cite{abatzoglou2014seasonal}.
%
However, how the co-variability between these two regions may evolve under climate change has not been systematically characterized, despite its importance for water management and climate resilience.
% \fix{ecosystem and agricultural planning across the region}.

Specifically, when the Bay-Delta experiences anomalously wet or dry water years, what spatial precipitation anomaly patterns does the broader West Coast realize, and is that relationship changing?
%
Using LOCA2-downscaled CMIP6 projections~\cite{pierce2014statistical,pierce2015improved,eyring2016overview} from 15 models under three emission scenarios~\cite{riahi2017shared} (SSP-2.45, 3.70, 5.85), we find that West Coast precipitation patterns conditioned on Bay-Delta wet and dry years are significantly departing from historical periods.
%
This departure indicates that the historical conditioning relationship between Bay-Delta and West Coast precipitation is itself changing, which we term ``decoherence''. We find that decoherence strengthens under higher emissions and over longer time horizons, most pronounced in spring months.

\subsection{Data}
\label{sec:data}
\subsubsection{LOCA2-Downscaled Climate Projections}
We analyze precipitation projections from 15 CMIP6 global climate models~\cite{eyring2016overview} (Table~\ref{tab:models}) downscaled to 6-km resolution over the contiguous United States using the Localized Constructed Analogs (LOCA2) method~\cite{pierce2014statistical,pierce2015improved,pierce2023future}.
%
LOCA2 provides monthly precipitation fields for historical (1950-2014) and future (2015-2100) periods under three Shared Socioeconomic Pathways: SSP2-4.5, SSP3-7.0, and SSP5-8.5~\cite{riahi2017shared}.
%
We also consider the first and second halves of the future period separately (2015-2057 and 2058-2100) to examine how decoherence evolves over time.
%
We consider two regions: the Sacramento Bay-Delta Watershed, a critical hydrological hub in California, and the broader U.S. West Coast region, encompassing coastal and inland areas from Washington to California (Figure~\ref{fig:regions}).

\begin{table}
\centering
\begin{tabular}{lccccc}
\hline
GCM & Variant & Historical & SSP2-4.5 & SSP3-7.0 & SSP5-8.5 \\
\hline
ACCESS-CM2       & r1i1p1f1 & \checkmark & \checkmark & \checkmark & \checkmark \\
CESM2-LENS       & r1i1p1f1 & \checkmark & \checkmark & \checkmark & \checkmark \\
CNRM-ESM2-1      & r1i1p1f2 & \checkmark & --         & \checkmark & --         \\
EC-Earth3-Veg    & r1i1p1f1 & \checkmark & \checkmark & \checkmark & \checkmark \\
EC-Earth3        & r1i1p1f1 & \checkmark & \checkmark & \checkmark & \checkmark \\
FGOALS-g3        & r1i1p1f1 & \checkmark & \checkmark & \checkmark & \checkmark \\
GFDL-ESM4        & r1i1p1f1 & \checkmark & \checkmark & \checkmark & \checkmark \\
HadGEM3-GC31-LL  & r1i1p1f3 & \checkmark & \checkmark & --         & \checkmark \\
INM-CM5-0        & r1i1p1f1 & \checkmark & \checkmark & \checkmark & \checkmark \\
IPSL-CM6A-LR     & r1i1p1f1 & \checkmark & \checkmark & \checkmark & \checkmark \\
KACE-1-0-G       & r1i1p1f1 & \checkmark & \checkmark & \checkmark & \checkmark \\
MIROC6           & r1i1p1f1 & \checkmark & \checkmark & \checkmark & \checkmark \\
MPI-ESM1-2-HR    & r1i1p1f1 & \checkmark & \checkmark & \checkmark & \checkmark \\
MRI-ESM2-0       & r1i1p1f1 & \checkmark & \checkmark & \checkmark & \checkmark \\
TaiESM1          & r1i1p1f1 & \checkmark & \checkmark & \checkmark & --         \\
\hline
\end{tabular}
\caption{Global climate models (GCMs) used in this study and their availability across SSP scenarios. Variant labels follow CMIP6 convention; absent combinations are marked with a dash and were unavailable at the time of the study.}
\label{tab:models}
\end{table}

\subsubsection{Data Preprocessing and Anomaly Calculation}
\begin{figure}
    \centering
    \includegraphics[width=0.75\textwidth]{west_coast_boundaries.pdf}
    \caption{Study region showing the Sacramento–Bay Delta Watershed (red) and the broader U.S. West Coast analysis domain (blue) used for examining precipitation co-variability.}
    \label{fig:regions}
\end{figure}

To isolate distributional changes from model biases and inter-model spread, we compute precipitation anomalies at each grid cell.
%
Here, we define the anomaly as the difference between the modeled precipitation, relative to the monthly climatology averaged across all 15 GCM models over the historical period (1950-2014).
%
This ensures our analysis detects shifts in spatial pattern distributions rather than absolute magnitude differences between models.

As in the original ClimateSOM, ensuring a pseudo-smooth set of SOM nodes that capture the full range of variability in the dataset is critical for characterizing the distribution of spatial patterns.
%
To this end, each monthly anomaly field is normalized by dividing by the standard deviation computed across all spatial locations, all SSP scenarios, all years, and all models for that specific month.
%
This normalization step is crucial for the SOM to learn meaningful patterns across months with different variability scales, preventing it from being dominated by months with higher variance.
%
Finally, we apply a signed-logarithm transformation to ensure negative and positive anomalies are equally represented in the SOM space.
%
Without transformation, the heavy positive tail of precipitation anomalies produces input vectors with extreme magnitudes that disproportionately influence SOM weight updates, causing the learned map to allocate excess capacity to representing rare extreme-wet states.
%
The signed-log compresses this asymmetry while preserving sign.

In all, for any precipitation field $p$ at month $m$ we compute the transformed anomaly $p_m'$:
\begin{linenomath*}
\begin{equation*}
    p_m' = \text{sign}(a_m) \cdot \log(|a_m| + 1), \quad a_m = \frac{p - \bar{p}_m}{\sigma_m}
\end{equation*}
\end{linenomath*}
where $\bar{p}_m$ is the historical monthly climatology (1950-2014) averaged across all 15 models for month $m$, and $\sigma_m$ is a scalar standard deviation computed across all spatial locations, SSP scenarios, years, and models for month $m$.
%
All transformed fields $p_m'$ across models, months, and time periods are pooled to train the SOM, ensuring it captures the full range of spatial patterns in the dataset.


\subsection{Methods}
\label{subsec:agu-methods}
\subsubsection{ClimateSOM Framework}
% We make use of the ClimateSOM framework with some adjustments.
% ClimateSOM~\cite{kawakami2025climatesom} was originally presented as a visual analysis tool that represents ensembles of spatial climate realizations as distributions in a two-dimensional space.
% %
% Here, we extend this framework to analyze sets of spatial realizations across different climate periods and scenarios.
%
% A self-organizing map (SOM) is trained on the transformed precipitation anomaly fields $p_m'$ (Section~\ref{sec:data}), producing a grid where each node represents a characteristic spatial pattern.
% %
% Given this SOM for any transformed precipitation field, we identify its best matching unit (BMU) in the SOM and map that node to a 2D position.
% %
% Aggregating these BMU positions across a set of realizations yields a distribution representing that set's spatial pattern composition.
% %
% Formally, for a set with $n$ transformed spatial fields $F=[f_1, f_2, \ldots, f_n]$, we obtain:
% \begin{equation}
% S=\{s_1,s_2,\dots,s_n\}, \,\,\, s_i=\text{proj}(\text{BMU}(f_i))
% \label{eq:climatesom_projection}
% \end{equation}
% where $\text{BMU}(f_i)$ is the best matching SOM node for transformed field $f_i$, and $\text{proj}(\cdot)$ maps the node to 2D coordinates.
% %
% In this work, we use a $30\times 30$ SOM grid, following the configuration established in the original ClimateSOM work~\cite{kawakami2025climatesom}. 
% \subsubsection{SOM Projection}
As identified, the original 2D SOM grid often contains distortions that we wish to mitigate.
%
The $\text{proj}(\cdot)$ function maps each SOM node to 2D coordinates that serve as the ground space for our EMD computation.
%
In this work, we use UMAP~\cite{mcinnes2018umap}, departing from the original ClimateSOM implementation, which used Minimum-Distortion Embedding (MDE)~\cite{agrawal2021minimum}.
%
MDE is a general framework that generalizes UMAP and t-SNE as special cases, and the additional anchoring constraints it exposes are not needed here, so we use UMAP as a simpler instance of the same family.

Local fidelity is the appropriate property for our application because EMD aggregates the cost of moving mass over short hops, and the climate-change signal we test takes the form of many such local redistributions in pattern composition rather than wholesale traversal across the pattern space.
%
Global distortions introduced by UMAP are therefore of limited concern, as long-range transports are less prevalent in our dataset. 
%
Instead, individual displacements are local, and their accumulation across realizations is what we report as decoherence.
%
The UMAP projection is computed once on the full SOM codebook, inspected for validity, and held fixed across all comparisons and permutations, so the test statistic depends only on group assignments and not on embedding stochasticity.

\subsubsection{EMD-Based Hypothesis Testing}
We extend ClimateSOM with a permutation test using Earth Mover's Distance (EMD) to formally compare distributions of spatial climate patterns between different sets of realizations.
% \shortnote{What if there's a different number?}
%
Recall that the SOM is trained once on the full dataset (all historical and future fields, all models and scenarios), ensuring a consistent pattern space for all comparisons.
%
Given two sets of spatial realizations with BMU distributions $S_1 = \{s_1^{(1)}, s_2^{(1)}, \ldots, s_{n_1}^{(1)}\}$ and $S_2 = \{s_1^{(2)}, s_2^{(2)}, \ldots, s_{n_2}^{(2)}\}$ in the UMAP-reduced ClimateSOM space, we compute the observed EMD:
\begin{equation}
d_{\text{obs}} = \text{EMD}(S_1, S_2)
\end{equation}
computed exactly via the network simplex solver in POT~\cite{flamary2021pot}. 
%
Here we use Euclidean distance between coordinates.
%
Additionally, we note that we solve the balanced optimal transport problem, which assumes equal mass in the compared distributions.
%
To satisfy this assumption, we equally distribute mass across all points in any set $S_i$ such that for any $S_i$ the total mass sums to 1.
%
This allows us to test distributions of arbitrary size under our framework.
%
The UMAP projection is computed once on the full SOM pattern space and held fixed across all permutations, ensuring the test statistic depends only on group assignments. 
% As an additional validation, we compare our EMD-based results against a standard two-sample t-test on the BMU coordinates and found qualitatively consistent significance patterns.

To assess statistical significance, we test the null hypothesis $H_0$ -- that the two distributions are drawn from the same underlying pattern distribution.
%
We perform a permutation test by pooling $S_1 \cup S_2$, randomly shuffling group labels, and recomputing EMD for each permutation.
%
Let $d_{\text{perm}}^{(k)}$ denote the EMD for the $k$-th permutation.
%
The p-value is calculated as:
\begin{equation}
p = \frac{1}{K} \sum_{k=1}^{K} I(d_{\text{perm}}^{(k)} \geq d_{\text{obs}})
\end{equation}
where $I(\cdot)$ is the indicator function.
%
We empirically find that $K=5000$ provides stable estimates and employ that here.
%
Lastly, we use $\alpha = 0.05$ as the significance threshold.
%
The validity of this permutation test rests on \emph{exchangeability} under $H_0$: relabeling realizations must be equally likely when the two sets are drawn from the same pattern distribution. We obtain this by treating each (model, water year) realization as an independent draw from the population of plausible climate states under its scenario. We acknowledge this independence assumption as a limitation, since realizations are correlated both across models, which share components and physics, and across neighboring water years; the resulting reduction in effective sample size may render our p-values mildly anti-conservative.

This approach detects both compositional reorganization (changes in which spatial patterns are realized) and pattern intensification, with one caveat.
%
The signed-log transformation compresses magnitude, so BMU assignment is pattern-dominated; intensified patterns shift to a more extreme neighboring node when one exists.
%
At the wettest and driest SOM extremes, however, no more-extreme nodes exist, so intensification of those patterns remains invisible to the test.
%
Significant decoherence therefore indicates a shift in the spatial pattern composition associated with a given regional regime, while non-significance does not strictly rule out intensification confined to the SOM extremes.

\subsubsection{Application to West Coast Precipitation Co-variability}

We apply this framework to test whether the relationship between Bay-Delta precipitation regime and West Coast spatial precipitation patterns may reorganize under climate change.
%
Specifically, we examine whether the spatial patterns realized across the broader West Coast during anomalously wet or dry Bay-Delta water years differ systematically between historical and future periods.

We classify each water year as wet or dry based on Bay-Delta accumulated precipitation, applying the same historical thresholds (defined below) to both historical and future periods so that wet/dry labels are consistent across all comparisons.
% %
% For each water year, we calculate the total accumulated precipitation in the Bay-Delta region.
%
Wet years are defined as those exceeding the 75th percentile of the historical distribution (1950-2014, pooled across all 15 models), and dry years fall below the 25th percentile.
%
For months within wet water years, we extract the corresponding West Coast spatial precipitation patterns from the historical period and compare them to West Coast patterns from months within wet water years in the future period.
%
All 15 climate models are pooled together in each comparison.
%
We repeat this procedure separately for dry water years.
%
Comparisons are conducted across October through May, three SSP scenarios (SSP2-4.5, SSP3-7.0, SSP5-8.5), and two future time periods (early: 2015-2057, late: 2058-2100).
%
% A significant p-value indicates that the West Coast precipitation response to a given Bay-Delta regime has fundamentally changed between historical and future climates.
% \subsection{Comparison and verification with standard analysis approaches}
% To assess the validity and robustness of our EMD-based hypothesis testing result, we compare our results

\subsection{Results}
\label{sec:results}
\begin{figure}
    \centering
    \includegraphics[width=0.9\textwidth]{som_grid_7x7.pdf}
    \caption{An evenly selected 7$\times$7 grid of SOM nodes from the trained ClimateSOM, showing the spatial precipitation anomaly patterns associated with each node. Values are in the log-standardized anomaly space used for training as  described in Section~\ref{subsec:agu-methods}. Broadly the SOM extracts two axes of change: (1) the North-South axes of wetness and dryness and (2) overall wetness and dryness across the entire West Coast. Each SOM node is projected via
    UMAP~\cite{mcinnes2018umap} and its projected position is used to represent any precipitation anomaly pattern in the dataset.
    \label{fig:som_grid}}
\end{figure}
% \subsection{Changes in Bay-Delta Wet and Dry Year Frequency}
%
% Before examining distributional changes in spatial patterns, we first test whether the frequency of Bay-Delta wet and dry water years is itself changing under climate scenarios.
% %
% Table~\ref{tab:wy_frequencies} shows the proportion of water years classified as dry (below 25th percentile) or wet (above 75th percentile) based on Bay-Delta accumulated precipitation across all 13 models for each scenario and time period.
% %
% Fisher's exact test reveals significant shifts in these proportions under higher emissions and in the late-century period.
% %
% SSP2-4.5 and SSP3-7.0 in the second half of the century show both a decrease in dry year frequency and a significant increase in wet year frequency, with wet years comprising 35.89\% and 32.75\% of water years respectively compared to the historical 25.00\%.
% %
% However, most notably, such changes do not occur under the highest emissions scenario (SSP5-8.5), where dry and wet year frequencies remain statistically indistinguishable from historical proportions in both future periods.
% %
% \shortnote{Whats this about?}
%
%
% \begin{table}
% \centering
% \caption{Frequency of Bay-Delta dry and wet water years across scenarios and time periods. Bold values differ significantly from historical proportions (Fisher's exact test, $p < 0.05$).}
% \label{tab:wy_frequencies}
% \begin{tabular}{c c c c}
% \hline
% Scenario & Period & Dry Model-Years (\%) & Wet Model-Years (\%) \\
% \hline
% Historical & 1950--2014 & 25.10 & 25.00 \\
% \hline
% SSP2-4.5 & Early (2015--2057) & 21.59 & 26.79 \\
%          & Late (2058--2100) & \textbf{20.21} & \textbf{35.89} \\
% \hline
% SSP3-7.0 & Early (2015--2057) & 25.16 & 26.62 \\
%          & Late (2058--2100) & \textbf{23.52} & \textbf{32.75} \\
% \hline
% SSP5-8.5 & Early (2015--2057) & 24.13 & 27.10 \\
%          & Late (2058--2100) & 24.58 & 29.46 \\
% \hline
% \multicolumn{4}{l}{\footnotesize Dry: $<$25th percentile; Wet: $\geq$75th percentile of historical distribution}
% \end{tabular}
% \end{table}

\subsubsection{Historical Model Evaluation Against Observations}

Before examining future changes, we evaluate whether the CMIP6 models capture observed historical co-variability patterns by comparing each model's historical period (1950-2014) against Livneh gridded observational data (1950-2014)~\cite{pierce2021extreme}, both as an (1) validation of each GCM to adequately \emph{replicate} historical patterns and (2) as a validation of our methodology.
%
We apply the same EMD-based permutation test to compare West Coast precipitation pattern distributions conditioned on Bay-Delta wet and dry years between Livneh observations and each GCM's historical simulation.
%
Figure~\ref{fig:livneh_unsplit_pvalues} shows p-values for these comparisons across months and conditioning states.
%
Most models show no significant decoherence from observations for most months, indicating that the models generally capture the historical co-variability structure between Bay-Delta regime and West Coast spatial patterns.
%
This provides confidence that detected future changes represent genuine reorganization signals rather than model artifacts.
%
% \begin{figure}
%     \centering
%     \includegraphics[width=\textwidth]{livneh_pvalue_table.pdf}
%     \caption{P-values from EMD permutation tests comparing Livneh observed precipitation patterns (1950-2011) against each CMIP6 model's historical simulation (1950-2014). Non-significant p-values indicate models that show no significant signs of deviations.}
%     \label{fig:livneh_pvalues}
% \end{figure}

\begin{figure}
    \centering
    \includegraphics[width=\textwidth]{livneh_unsplit_pvalue_table.pdf}
    \caption{P-values from EMD permutation tests comparing Livneh unsplit observed precipitation patterns (1950-2014) against each CMIP6 model's historical simulation (1950-2014). Non-significant p-values indicate models that show no significant signs of deviations.}
    \label{fig:livneh_unsplit_pvalues}
\end{figure}

\subsubsection{Welch t-test based hypothesis testing results}\label{ssec:t_test_results}
To assess the validity of our EMD-based decoherence test, we compare its results against a standard parametric approach using Welch's t-test on mean West Coast precipitation anomalies.
%
Welch's variant does not assume equal variances, an appropriate choice given the differing sample sizes and variance characteristics across periods and scenarios.

The two tests address related but distinct questions.
%
The t-test detects shifts in the mean magnitude of anomalies between periods, whereas the EMD test detects shifts in the distribution of spatial patterns themselves.
%
Agreement between the two provides converging evidence of distributional change; disagreement would highlight cases where pattern composition shifts without a corresponding shift in mean intensity.

We compute the t-test for each combination of month, scenario, and conditioning state, comparing all grid-point anomalies in the West Coast domain between historical and future periods.
%
Results appear alongside the EMD results in Figure~\ref{fig:combined_pvalues}.
%
Broadly, the two tests show consistent significance patterns, supporting the validity of our EMD-based approach.
%
Across all 96 test cells (3 SSP scenarios $\times$ 2 future half-periods $\times$ 2 Bay-Delta conditioning states $\times$ 8 months), the two tests reach the same significance verdict at $p < 0.05$ in 71 cases (74\%), as summarized in Table~\ref{tab:emd_welch_agreement}.
%
Critically, the disagreements are highly asymmetric: in 22 of the 25 discordant cells the EMD test detects a significant distributional shift where the Welch t-test does not, whereas the reverse occurs in only 3 cells.
%
This asymmetry is consistent with our central premise, that the spatial composition of precipitation patterns can reorganize without a detectable change in mean intensity, and that such reorganization is visible to a distributional test yet invisible to a mean-based one.
%
\begin{table}[t]
    \centering
    \begin{tabular}{lccc}
        \toprule
        & \multicolumn{2}{c}{Welch t-test} & \\
        \cmidrule(lr){2-3}
        & Significant & Not significant & Total \\
        \midrule
        EMD significant     & 19 & 22 & 41 \\
        EMD not significant & 3  & 52 & 55 \\
        \midrule
        Total               & 22 & 74 & 96 \\
        \bottomrule
    \end{tabular}
    \caption{Agreement between the EMD-based permutation test and the Welch t-test across all 96 test cells (3 SSP scenarios $\times$ 2 future half-periods $\times$ 2 Bay-Delta conditioning states $\times$ 8 months), with each cell classified as significant or not at $p < 0.05$. The tests agree in 71 of 96 cells (74\%): 19 jointly significant and 52 jointly non-significant. Of the 25 disagreements, 22 are cells where the EMD test detects a distributional shift that the mean-based Welch test misses, and only 3 are the reverse.}
    \label{tab:emd_welch_agreement}
\end{table}
%
The following sections focus on the EMD-based analysis.
% \begin{figure}
%     \centering
%     \includegraphics[width=\textwidth]{welch_pvalue_table.pdf}
%     \caption{P-values from Welch t-tests comparing mean West Coast precipitation anomalies conditioned on Bay-Delta wet and dry years between historical (1950-2014) and future periods (2015-2057, 2058-2100) across three SSP scenarios }
%     \label{fig:t_test_pvalues}
% \end{figure}

\subsubsection{EMD-based hypothesis testing results}
In the following section, we evaluate the future vs historical West Coast precipitation distributional differences using our ClimateSOM extended EMD-based method, instead of considering the mean of all precipitation anomaly patterns as was performed above ~\ref{ssec:t_test_results}.
% 
Again the precipitation patterns are conditioned upon Bay-Delta watershed precipitation being wet or dry.
%
Our permutation tests reveal significant decoherence in West Coast precipitation pattern distributions conditioned on Bay-Delta regime from historical to future periods (Figure~\ref{fig:combined_pvalues}).
%
The magnitude and timing of decoherence depend systematically on emission scenario, time horizon, conditioning state (wet vs. dry), and month.

% \begin{figure}
%     \centering
%     \includegraphics[width=\textwidth]{pvalue_table.pdf}
%     \caption{P-values from Earth Mover's Distance permutation tests (1,000 permutations) comparing distributions of West Coast spatial precipitation patterns conditioned on Bay-Delta wet and dry years between historical (1950-2014) and future periods. Results span 13 CMIP6 models aggregated, three SSP scenarios (rows), two future time periods (early: 2015-2057, late: 2058-2100), Bay-Delta wet and dry conditioning states, and eight months (Oct-May). Significant decoherence ($p < 0.05$) strengthens with higher emissions, later time horizons, and varies systematically by month and conditioning state.}
%     \label{fig:pvalues}
% \end{figure}

\begin{sidewaysfigure}
    \centering
    \includegraphics[width=\textwidth]{combined_pvalue_table.pdf}
    \caption{Combined p-values from Welch t-tests and EMD-based permutation tests comparing historical and future distributions of West Coast precipitation patterns conditioned on Bay-Delta wet and dry years. }
    \label{fig:combined_pvalues}
\end{sidewaysfigure}

Decoherence strengthens from early-century to late-century periods.
%
In the early future period (2015-2057), Bay-Delta dry years show emerging decoherence primarily in October and April-May under higher emissions (SSP5-8.5: $p = 0.000$ in October, $p = 0.017$ in April), while wet years show weaker signals with only scattered significant months.
%
Table~\ref{tab:emd_sig_counts} summarizes this progression by tallying, for each scenario and conditioning regime, the number of months (out of eight, October--May) that exhibit significant distributional divergence ($p < 0.05$) under the EMD test.
%
The tally of significant months more than doubles from the early to the late period, from 13 to 28 summed across all scenarios and conditioning states, confirming that decoherence broadens across the wet season as the century progresses rather than remaining confined to a handful of months.
%
The signal also grows monotonically with emissions within nearly every period-regime block and is consistently stronger in dry-year than wet-year conditioning during the first half, with wet years catching up only in the late century.

\begin{table}[t]
    \centering
    \begin{tabular}{lcccc}
        \toprule
        & \multicolumn{2}{c}{First-Half (2015--2057)} & \multicolumn{2}{c}{Second-Half (2058--2100)} \\
        \cmidrule(lr){2-3} \cmidrule(lr){4-5}
        Bay-Delta Conditioning & Dry & Wet & Dry & Wet \\
        \midrule
        SSP2-4.5 & 3 & 1 & 5  & 4  \\
        SSP3-7.0 & 2 & 1 & 4  & 4  \\
        SSP5-8.5 & 4 & 2 & 6  & 5  \\
        \bottomrule
    \end{tabular}
    \caption{Number of months (out of eight, October--May) with significant EMD-based distributional divergence ($p < 0.05$) between historical and future West Coast precipitation pattern distributions, by SSP scenario, future half-period, and Bay-Delta conditioning regime.}
    \label{tab:emd_sig_counts}
\end{table}

%
The late-century period (2058-2100) reveals widespread decoherence across both conditioning states.
%
For Bay-Delta dry years, fall and early winter months (October-January) show strong significance across all SSP scenarios, particularly under SSP3-7.0 and SSP5-8.5, and spring months (April-May) maintain significance across scenarios.

Bay-Delta wet years in the late period exhibit the strongest reorganization signals.
%
Under SSP5-8.5, late winter through spring (March-May) shows highly significant decoherence ($p = 0.000$ for all three months), with additional shifts in November ($p = 0.003$) and February ($p = 0.019$).
%
SSP3-7.0 shows similar spring decoherence (April-May: $p = 0.000$) and significant shifts in November ($p = 0.007$) and February ($p = 0.000$).
%
Additionally, SSP2-4.5 shows significant late-century wet-year decoherence in March-May.

Winter core months (December-January) generally show weaker signals, particularly for wet years.
%
This may indicate that mid-winter co-variability is undergoing intensification of existing spatial patterns rather than compositional reorganization; such changes would not be detected by our BMU distribution comparison, as intensified patterns would still map to the same SOM nodes.
%
The asymmetry between wet and dry year responses, especially pronounced in the late century, indicates that the Bay-Delta relationship with West Coast patterns is reorganizing in regime-dependent ways rather than through simple uniform shifts.
%
These results demonstrate that historical co-variability structures are fundamentally breaking down under climate change, with the rate and character of reorganization varying systematically by emission pathway, time horizon, and regional conditioning state.

% \begin{figure}
%     \centering
%     \includegraphics[width=\textwidth]{median_pvalue.png}
%     \caption{Median p-values across MIROC6, EC-Earth3-Veg, KACE-1-0-G, ACCESS-CM2, IPSL-CM6A-LR, INM-CM5-0 comparing pooled historical and future ssp370}
%
% \end{figure}


% is not enough to detect significant differences. I think the best way to show this is to just show the distribution of p-values across models for each cell in Figure 4, but even then I think it will be hard to interpret. Let me know if you want me to try that or if you have other ideas.}

\subsubsection{Visualizing the Overall Spatial Character of Reorganization}\label{ssec:cumulative_change}

Given the significant decoherence detected in the previous section, we now visualize \textit{what} is changing spatially.
%
To this end, we compute cumulative change fields where each SOM node's spatial pattern is weighted by the shift in its occupancy frequency.
%
For each comparison between distributions $S_1$ (future) and $S_2$ (historical), we calculate the proportion of samples matching each SOM node $i$ in both distributions. 
\begin{align*}
    p_{1,i} &= \text{count}(S_1 \text{ matches to node } i)/|S_1| \\
    p_{2,i} &= \text{count}(S_2 \text{ matches to node } i)/|S_2|
\end{align*}
The cumulative change field is then
\begin{linenomath*}
\begin{equation*}
    \text{Cumulative Change} = \sum_{i} (p_{1,i} - p_{2,i}) \cdot W_i
\end{equation*}
\end{linenomath*}
where $W_i$ is the spatial precipitation pattern associated with SOM node $i$.
%
Positive values indicate locations whose expected precipitation under the given conditioning is higher in the future than historically; negative values, lower.
%
We note that the field is a relative change between historical and future distributions and does not characterize the absolute wetness or dryness of locations as projected in the CMIP6 catalog.


\begin{sidewaysfigure}
    \centering
    \begin{subfigure}{0.49\textwidth}
        \centering
        \includegraphics[width=\textwidth]{change_mosaic_grid_first_dry.pdf}
        \caption{Dry Bay-Delta years, early future (2015--2057)}
        \label{fig:change_mosaic_grid_first_dry}
    \end{subfigure}
    \hfill
    \begin{subfigure}{0.49\textwidth}
        \centering
        \includegraphics[width=\textwidth]{change_mosaic_grid_second_dry.pdf}
        \caption{Dry Bay-Delta years, late future (2058--2100)}
        \label{fig:change_mosaic_grid_second_dry}
    \end{subfigure}

    \vspace{0.75em}

    \begin{subfigure}{0.49\textwidth}
        \centering
        \includegraphics[width=\textwidth]{change_mosaic_grid_first_wet.pdf}
        \caption{Wet Bay-Delta years, early future (2015--2057)}
        \label{fig:change_mosaic_grid_first_wet}
    \end{subfigure}
    \hfill
    \begin{subfigure}{0.49\textwidth}
        \centering
        \includegraphics[width=\textwidth]{change_mosaic_grid_second_wet.pdf}
        \caption{Wet Bay-Delta years, late future (2058--2100)}
        \label{fig:change_mosaic_grid_second_wet}
    \end{subfigure}
    \caption{Cumulative change fields of West Coast precipitation relative to the historical period (1950--2014), aggregated across 15 CMIP6 models by month (columns) and SSP scenario (rows within each panel), conditioned on Bay-Delta dry years (top) and wet years (bottom) for the early (left) and late (right) halves of the future period. Positive (teal) values indicate locations wetter in the future than historically under the given conditioning, and negative (brown) values drier; bordered cells denote months with significant distributional divergence ($p < 0.05$).}
    \label{fig:change_mosaic_grid}
\end{sidewaysfigure}

%
%
%
%

\subsubsection{Shifting seasonal precipitation in the U.S. West Coast, conditioned on Bay-Delta wet and dry years}
\paragraph*{Bay-Delta Dry Years}
Bay-Delta dry years show the strongest decoherence in fall (October-November) and spring (April-May), particularly under higher emissions and in the late future period.
%
The cumulative change fields reveal that these changes are primarily driven by a shift of precipitation away from the Bay-Delta Watershed toward the Pacific Northwest, with the exception of May in both periods where Bay-Delta dryness co-occurs with an overall dry U.S. West Coast (Figures~\ref{fig:change_mosaic_grid_first_dry} and~\ref{fig:change_mosaic_grid_second_dry}).

\paragraph*{Bay-Delta Wet Years}
Bay-Delta wet years show modest decoherence in the early future period, but exhibit more systematic and stronger decoherence in the late future period, particularly under SSP5-8.5.
%
The most notable changes occur in spring months (MAM), especially May, where Bay-Delta wet water years associate with a shift toward drier conditions in California and the Sierra Nevada, and wetter conditions in the Pacific Northwest (Figure~\ref{fig:change_mosaic_grid_second_wet}).
%
% This suggests that even during anomalously wet water years in the Bay-Delta, spring months may experience relative dryness compensated by wetter conditions in other months of the water year \fix{this needs to be verified}.
%
Similar trends appear in November under SSP3-7.0 and SSP5-8.5 in the late future period.


\subsubsection{Within-year coherence of cumulative changes}
The cumulative change fields shown in Figures~\ref{fig:change_mosaic_grid_first_dry},~\ref{fig:change_mosaic_grid_second_dry},~\ref{fig:change_mosaic_grid_first_wet}, and~\ref{fig:change_mosaic_grid_second_wet} represent aggregate patterns computed across all years within each period-condition.
%
However, a critical question remains unanswered. 
%
Do these spatial patterns actually manifest together within individual water years, or do they emerge only when averaging across the entire period?
%
To address this, we test whether the cumulative change patterns are coherently expressed in specific water years at rates higher than expected by chance.

For each (month $m$, condition $c$), we first compute the historical conditional mean anomaly field $\bar{H}_{m,c}$ by averaging historical anomalies for that month and condition.
%
For each future month-model-year under condition $c$, we then compute the delta from this baseline:
\begin{equation*}
\delta_{m,y,k,c} = a_{m,y,k} - \bar{H}_{m,c}
\end{equation*}
This delta is directly comparable to the cumulative change field, since both are deltas relative to historical conditional behavior.
%
The subtraction also removes confounding from the historical conditional pattern (for example, the typical wet-year pattern that already existed historically) bleeding into the comparison.
%
We then calculate the Pearson correlation between each per-month delta $\delta_{m,y,k,c}$ and the corresponding cumulative change field, and average these monthly correlations across Oct-May for each model-year, yielding a single coherence metric per model-year.
%
A model-year with high mean correlation indicates that multiple months within that year exhibit shifts spatially aligned with the cumulative change field.

We use 0.3 as a threshold for identifying model-years with meaningful coherence.
%
While modest, this threshold provides a conservative indicator of systematic pattern expression.
%
We count the number of model-years exceeding this threshold within each period-condition as our test statistic.

To determine whether the observed counts are statistically significant, we perform another permutation test.
%
For each month within a period-condition, we randomly reassign all correlation values for that month across different year-model combinations.
%
This shuffling preserves the distribution of correlation values for each month but removes any temporal coherence within individual water years.
%
We recalculate model-year mean correlations from these shuffled values and count how many exceed 0.3, repeating this process to build a null distribution.
%
Under the null hypothesis that the cumulative changes are purely an aggregate artifact with no systematic expression in individual years, the observed count would be consistent with the null distribution.
%
Observed counts significantly exceeding the null distribution indicate that specific water years do indeed exhibit the cumulative change patterns coherently across multiple months.
% we perform a permutation test on the count of year-model groups with mean correlation $> 0.3$ between the change fields for each period-condition.
\begin{table}
    \centering
\begin{tabular}{l r r r r}
\hline
Period-Condition & Observed Count ($>0.3$)  & Null Mean & Null SD & p-value \\
\hline
First - Dry & 26  & 19.83 & 3.87 & 0.0778 \\
First - Wet & 51  & 46.48 & 5.15 & 0.2190 \\
Second - Dry & 93  & 76.12 & 5.62 & \textbf{0.0018} \\
Second - Wet & 113 & 91.91 & 6.45 & \textbf{0.0012} \\
\hline
\end{tabular}
\caption{Results of permutation test for co-occurrence of cumulative change fields with monthly anomaly fields across model-years for each period-condition ($\alpha=0.05$). Significant p-values indicate that the observed count of model-years with mean correlation $>0.3$ is greater than expected under the null distribution generated by randomly reassigning correlation values within each month across year-model combinations.}
\label{tab:per_period_pvals}
\end{table}
%
We show the results in Table~\ref{tab:per_period_pvals}.
%
The results demonstrate that the co-occurrence of cumulative changes is not merely an artifact of aggregating across entire periods, but rather reflects patterns coherently expressed in individual model-years at rates significantly higher than expected by chance.
%
This is particularly evident in the second half of the century for both dry and wet Bay-Delta years, where observed counts substantially exceed null expectations.
%
% Figure~\ref{fig:annual_cooccurrence} shows the number of model-years with mean correlation $>0.3$ for each year within each period-condition, illustrating the temporal evolution of within-year coherence.

% \begin{figure}
%     \centering
%     \includegraphics[width=0.7\textwidth]{annual_cooccurence.png}
%     \caption{Annual count of model-years with mean correlation $>0.3$ between monthly anomaly fields and cumulative change fields for each period-condition. The increase in counts in the second half of the century visually supports the significant co-occurrence results from the permutation test.}
%     \label{fig:annual_cooccurrence}
% \end{figure}

% \subsection{Discussion}
% \label{sec:discussion}
% \shortnote{To add: Implications, Limitations}

\subsection{Continued Work}
\label{subsec:agu-conclusion}

This section forms the last of three arms in our extension of ClimateSOM.
%
First, we extended the workflow with a statistical framework, an EMD-based permutation test that determines whether the distribution of conditioned spatial precipitation patterns has shifted under a changing climate.
%
Second, we established the validity of this framework by validating it against a traditional approach, Welch's t-test on mean precipitation, and found the two in broad agreement.
%
Third, in this final arm, we develop how to visualize the change that the test detects; this work is ongoing, and we describe it here to outline where we are currently taking the extension.

Our aim is to visualize the spatial reorganization that our EMD test has detected.
%
The cumulative change fields above give a first, aggregate view of this reorganization, and we are now developing methods to resolve that view further, decomposing the aggregate change into the specific spatial pattern transitions that compose it, so that the detected reorganization can be visualized as a set of interpretable shifts between precipitation patterns rather than a single blended field.

\subsubsection{Decomposing the Spatial Reorganization}\label{ssec:decomposing_reorganization}
The cumulative change field $\sum_i (p_{1,i} - p_{2,i}) W_i$ defined in Section~\ref{ssec:cumulative_change} aggregates contributions from many SOM nodes and therefore does not isolate the specific pattern transitions responsible for the observed reorganization.
%
To resolve this, we are developing a clustering-based decomposition (Algorithm~\ref{alg:decomposition}) that couples optimal transport with density-based clustering of SOM-projected samples in the ClimateSOM space (Figure~\ref{fig:agu_som_space}), quantifying mass redistribution between historical and future periods at the level of pattern clusters.
%
We select HDBSCAN over $k$-means and hierarchical clustering because it does not require the number of clusters to be specified \textit{a priori} and accommodates clusters of varying shapes and densities. Points classified by HDBSCAN as noise are excluded from the decomposition.

\begin{figure}[t]
    \centering
    \includegraphics[width=\textwidth]{AGU_SOM_Space.pdf}
    \caption{The two-dimensional ClimateSOM space over which the decomposition operates. Each point is a SOM node placed at its projected position, and the callout shows the log-standardized precipitation anomaly pattern associated with one node. Every sample is located in this space by its best-matching unit, so a set of samples becomes a density over the space that the decomposition then clusters and transports.}
    \label{fig:agu_som_space}
\end{figure}

The reduced transport plan $\tilde{P}$ quantifies the mass redistributed between each pair of pattern clusters; however, redistributed mass alone is an insufficient measure of the magnitude of spatial change.
%
A large transfer between two similar clusters constitutes a minor reorganization, whereas a smaller transfer between dissimilar clusters may contribute substantially to the cumulative change field.
%
To rank transitions by their contribution to the overall reorganization, we define an \textit{importance score} for each cluster pair as the product of (1) the redistributed mass $\tilde{P}_{ab}$ and (2) the mean absolute error (MAE) between the cluster centroids in the SOM prototype space.
%
The MAE between centroids in the SOM prototype space directly measures the dissimilarity between the underlying spatial precipitation patterns. High importance scores then identify transitions in which substantial mass is exchanged between markedly distinct patterns.

Although unnormalized importance scores have different totals across comparisons, reflecting different magnitudes of total reorganization, expressing each as a percentage of its comparison's total yields a compositional measure: the fraction of that comparison's spatial reorganization attributable to a given cluster transition.
%
This normalized score lets us rank the cluster transitions driving the observed changes and compare their contributions across comparisons, even when the underlying comparisons differ in scale or sample size. In practice, we examine the top-ranked cluster transitions per comparison to identify the dominant pattern shifts.
%
Figure~\ref{fig:agu_decomposition} illustrates the decomposition for an example comparison: historical versus SSP5-8.5 March precipitation in the late-century period (2058--2100), conditioned on Bay-Delta wet years.
%
We select this comparison because the EMD test flags it as a significant distributional shift (Figure~\ref{fig:combined_pvalues}, $p = 0.001$), so a genuine reorganization is known to be present for the decomposition to resolve.
%
The before and after distributions are each partitioned into clusters, and the optimal transport plan redistributes probability mass between them to isolate the dominant cluster transitions.
%
Figure~\ref{fig:agu_centroid_weights} then renders, for the five highest-scoring transitions, the source and target cluster centroids as spatial precipitation patterns, making the specific reorganization each transition represents directly legible.

Conditioned on Bay-Delta wet March years in the late century under SSP5-8.5, the five highest-scoring transitions share a consistent theme: a loss of precipitation in the south giving way to gains in the north.
%
The top transition (Flow~1, importance score $17.35\%$) relaxes an extreme-dry Sierra Nevada and Northern California source pattern toward a more moderate dryness.
%
Flow~2 ($14.49\%$) loses a wet Southern California pattern to a slightly dry Southern California and Sierra state.
%
Flow~3 ($12.49\%$) introduces a new wet Northern California and Sierra pattern.
%
No single transition dominates, but instead the reorganization spreads across several comparably-weighted flows that together shift wet-year March precipitation away from Southern California and toward Northern California and the Pacific Northwest.
%
\begin{algorithm}[h!]
\caption{Clustering-based decomposition of distributional reorganization}
\label{alg:decomposition}
\begin{algorithmic}[1]
\Require Samples $S_1$ with $|S_1| = n$ and samples $S_2$ with $|S_2| = m$, projected onto the SOM space
\Ensure Reduced transport plan $\tilde{P}$ and normalized importance scores $\tilde{I}$ summarizing cluster-level reorganization between $S_1$ and $S_2$
\State Compute the optimal transport plan $P \in \mathbb{R}^{n \times m}$ between $S_1$ and $S_2$ using their BMU locations in the projected SOM space, where $P_{ij}$ denotes the mass transported from sample $i \in S_1$ to sample $j \in S_2$. Use normalized mass if $n\neq m$.
\State Apply HDBSCAN to the projected SOM locations of $S_1$ and $S_2$ separately, yielding clusters $C_a$ ($a \in \{1, \dots, k_1\}$) and $D_b$ ($b \in \{1, \dots, k_2\}$); exclude points labeled as noise
\State Compute centroids $c_a$ for each $C_a$ and $d_b$ for each $D_b$ as the mean SOM prototype of their member nodes
\For{each cluster pair $(a, b)$}
    \State $\tilde{P}_{ab} \gets \sum_{i \in C_a} \sum_{j \in D_b} P_{ij}$
    \State $I_{ab} \gets \tilde{P}_{ab} \cdot \mathrm{MAE}(c_a, d_b)$ \Comment{importance score }
\EndFor
\State $Z \gets \sum_{a, b} I_{ab}$ \Comment{total importance across all cluster pairs}
\For{each cluster pair $(a, b)$}
    \State $\tilde{I}_{ab} \gets I_{ab} / Z$ \Comment{normalized importance score}
\EndFor
\State \Return $\tilde{P}, \tilde{I}$
\end{algorithmic}
\end{algorithm}

\begin{figure}[p]
    \centering
    \begin{subfigure}{0.49\textwidth}
        \centering
        \includegraphics[width=\textwidth]{AGUBeforeDist.png}
        \caption{Before (historical) distribution}
        \label{fig:agu_before_dist}
    \end{subfigure}
    \hfill
    \begin{subfigure}{0.49\textwidth}
        \centering
        \includegraphics[width=\textwidth]{AGUAfterDist.png}
        \caption{After (future) distribution}
        \label{fig:agu_after_dist}
    \end{subfigure}

    \vspace{0.75em}

    \begin{subfigure}{0.8\textwidth}
        \centering
        \includegraphics[width=\textwidth]{TransportPlan.png}
        \caption{Optimal transport plan between clusters}
        \label{fig:agu_transport_plan}
    \end{subfigure}
    \caption{Clustering-based decomposition of an example comparison: historical versus SSP5-8.5 March precipitation in the late-century period (2058--2100), conditioned on Bay-Delta wet years. The before (a) and after (b) sample distributions are each represented as a density over the ClimateSOM space and partitioned into clusters, labeled by their share of the probability mass, with markers at the cluster centroids. The optimal transport plan (c) redistributes mass from the before-clusters (circles) to the after-clusters (diamonds). The \emph{historical}
    distributions shown as the same density rendering and the \emph{future} distributions in the dotted outline. The labeled arrows give the dominant cluster transitions and its importance score $\tilde{I}$.}
    \label{fig:agu_decomposition}
\end{figure}

\begin{figure}[p]
    \centering
    \includegraphics[width=\textwidth]{centroid_weights.pdf}
    \caption{The five highest-scoring cluster transitions (flows) from the transport plan in Figure~\ref{fig:agu_transport_plan}, ranked by importance score. Each column is one transition; the top row shows the source (historical) cluster centroid and the bottom row the target (future) cluster centroid, both rendered as log-standardized precipitation anomaly patterns. The annotations report the redistributed mass, the MAE between the two centroids ($\mathrm{Dist}$), and the resulting normalized importance score. Rendering the centroids directly turns each ranked transition into an interpretable before-and-after pair of spatial precipitation patterns.}
    \label{fig:agu_centroid_weights}
\end{figure}

% \subsection{Conclusion}
\subsection{Conclusion}
\label{subsec:agu-real-conclusion}
This extension carried ClimateSOM from exploration to inference.
%
We equipped the workflow with an EMD-based permutation test that asks, with statistical rigor, whether the distribution of conditioned spatial precipitation patterns has shifted under a changing climate, and we validated it against the standard Welch t-test, finding broad agreement together with an informative asymmetry: the distributional test detects reorganizations of spatial pattern that a mean-based test did not detct.
%
Applied to the co-variability between the Sacramento Bay-Delta Watershed and U.S. West Coast precipitation, the test reveals a statistically significant, regime-dependent decoherence that strengthens with emissions and into the late century, and that is coherently expressed within individual water years rather than emerging only in the period average.
%
The clustering-based decomposition then begins to turn these detected shifts into interpretable spatial pattern transitions, as in the worked example where wet-year March precipitation is projected to redistribute from Southern California toward the north.

% We regard this as a successful sarlitmus test for ClimateSOM.
%
Confronted with a concrete, climatologically substantive question and with the established tools of the domain, the workflow delivers an inferential result those tools were not designed to provide, and renders that result as legible spatial patterns rather than a single summary statistic.
%
The decomposition remains a work in progress, and our broader aim is to fold this inferential and interpretive capability back into the interactive ClimateSOM workflow, so that hypothesis testing and pattern attribution become steps an analyst can take within the same visual environment used to explore the ensemble in the first place.


